% !TeX root = ../ba-lentz.tex

\documentclass[ba-lentz.tex]{subfile}

\begin{document}

{\let\clearpage\relax \chapter{Erwartete Ergebnisse: Verbesserung von saisonalen Vorhersagen und die Überbrückung von Beobachtungs-Vorhersage Lücken}}\label{results}

Durch das hybride Nachbearbeiten dynamischer Vorhersagen mithilfe von KI kann eine signifikante Verbesserung des Vorhersageskills der $KWB$ und damit auch des $SPEI$ erreicht werden.
Diese Erwartung zieht sich einerseits aus der statistischen Korrektur von Modellbiase und andererseits, dass die verwendeten PCNN nur auf dem afrikanischen Kontinent kalibriert werden für die Vorhersage der $KWB$.
Die ursprünglichen dynamischen Vorhersagen dagegen sind global auf die generelle Vorhersage von Temperatur und Niederschlag kalibriert. \\
Weiterhin soll versucht werden die Auflösungslücke zwischen Remote Sensing Daten und dynamischen Vorhersagen durch ein möglichen implementierten Downscaling-Mechanismus zumindest verkleinert werden.
Dabei ist die Funktion eines solchen Mechanismus aufgrund seiner Neuheit schwer einzuschätzen.
Da aber das verwendete PCNN großen Skill in dem ausfüllen von Klimadaten bewiesen hat \parencite{lentz_improving_2025} und Downscaling im Grunde genommen auch in einem Ausfüllen von neuen Lücken besteht, kann auch hier von einer zumindest einigermaßen realistischen Erhöhung der Vorhersagauflösung ausgegangen werden.
In wieweit eine solche höheraufgelöste Vorhersage tatsächlich auch regionale Dürreereignisse besser vorhersagen kann, ist dabei nicht vorher abschätzbar und muss im Arbeitsprozess untersucht werden.




\end{document}